% !TeX root = ../main.tex
\section{Cài Đặt và Phần Mềm}
\label{sec:impl}

Game được viết bằng Python~3 với thư viện Pygame~2. Mã nguồn mở tại:
\url{https://github.com/quyetops/Tactical-Formation-The-Board-Game}.

Mã nguồn được tổ chức thành các module độc lập: \texttt{core/} (engine
luật, hình học bàn cờ, trạng thái), \texttt{ai/} (bot Minimax,
Paranoid, chính sách), \texttt{ui/} (render Pygame, screens, widgets,
animation), và \texttt{tests/} (kiểm thử luật). Tài liệu kỹ thuật
\texttt{docs/formal\_model.md} mô tả mô hình toán học hình thức và
\texttt{docs/visual\_model.html} cung cấp companion trực quan.

\subsection{Ánh xạ mô hình hình thức sang mã nguồn}
\label{ssec:mapping}

Module \texttt{core/} hiện thực trực tiếp các định nghĩa ở
Phụ lục~\ref{app:formal}, giữ đúng nguyên tắc \emph{engine luật không phụ
thuộc giao diện hay AI}:
\begin{itemize}[noitemsep,topsep=2pt]
  \item Trạng thái~\eqref{eq:state} là một cấu trúc bất biến gồm cấu hình
        bàn $b$ (ánh xạ ô $\to$ màu) và phe đến lượt.
  \item Sinh nước đi cài đặt $\mathrm{Jumps}(x,b)$~\eqref{eq:moves}: với
        mỗi quân của phe đến lượt, duyệt tập hướng $D_{\kappa(x)}$ và lọc
        theo ba điều kiện hợp lệ~\eqref{eq:jump}.
  \item Phát hiện kẹp cài đặt phép quét tia~\eqref{eq:clamp}: từ ô đáp,
        theo mỗi hướng ăn, gom dãy quân địch liền nhau cho tới khi gặp
        quân mình (thành công), ô trống, hoặc biên bàn (thất bại).
  \item Áp dụng nước đi là hàm chuyển $T$~\eqref{eq:t1}--\eqref{eq:t3}:
        nhảy, lật theo hướng đã chọn, đổi lượt---thuần tất định, không
        sinh quân, không xoá quân (kiểm bất biến~\eqref{eq:conservation}
        bằng một assertion tổng quân $=28$ trong \texttt{tests/}).
\end{itemize}
Bot trong \texttt{ai/} chỉ gọi ba phép trên (sinh nước, áp dụng, đánh
giá): \texttt{minimax.py} cài Thuật toán~\ref{alg:minimax} cho chế độ 2
người, \texttt{paranoid.py} cài Thuật toán~\ref{alg:paranoid} cho chế độ
3--4 người. Cả hai tách biệt hoàn toàn với engine luật và có thể thay
bằng thuật toán khác mà không sửa \texttt{core/}.

\begin{lstlisting}[language=bash, caption={Cài đặt và chạy game.}]
python -m venv .venv
source .venv/bin/activate   # Linux / macOS
pip install pygame
python main.py
\end{lstlisting}

Phiên bản web chạy qua \texttt{pygbag}~\citep{pygbag}:
\begin{lstlisting}[language=bash, caption={Chạy trên trình duyệt.}]
pip install pygbag
pygbag main.py
# Mở trình duyệt tại http://localhost:8000
\end{lstlisting}
